\documentclass[10pt,a4paper]{amsart}
\usepackage[margin=19mm]{geometry}
\usepackage{amsmath,amssymb,mathtools}
\usepackage[scr=rsfs,cal=euler]{mathalfa}
\usepackage{xfrac}
\usepackage[unicode,hidelinks,bookmarks=false]{hyperref}
\newcommand{\ie}{i.e.\ }
\newcommand{\cf}{cf.\ }
\newcommand{\resp}{respectively\ }
\newcommand{\Subsection}{\subsection}
\providecommand{\nobreakdash}{\nobreak\hbox{-}}
\setlength{\emergencystretch}{2em}
\multlinegap0pt
\usepackage{amssymb}
\usepackage{tikz}
\newtheorem{lemma}{Lemma}
\title{Order symmetry and orthogonality of trajectories in discrete interval exchange transformations}
\author{S\'ebastien S. Ferenczi, Luca Q. Zamboni}
\date{Source: arXiv:2607.03785v1 (CC BY 4.0)}
\begin{document}
\maketitle

\noindent\emph{Source and licence.} Order symmetry and orthogonality of trajectories in discrete interval exchange transformations. Version arXiv:2607.03785v1 (CC BY 4.0), distributed by arXiv under the Creative Commons Attribution 4.0 International licence, http://creativecommons.org/licenses/by/4.0/, as recorded in the arXiv licence metadata for this exact version and shown on the abstract page https://arxiv.org/abs/2607.03785v1. Full licence notice: \texttt{licenses/arxiv-2607.03785-CC-BY-4.0.txt}.\par\smallskip

\noindent\textit{Editorial scope.} Counted unit: the article's lemma code (source0.tex lines 474--481). The article's complete original proof of it is delivered with it (source0.tex lines 483--510, one \texttt{proof} environment of 28 lines). No mathematics is written, repaired or re-derived: the statement and the proof are the article's own lines, in the article's own notation, and no retained line is substituted or re-flowed.  No further source-local context is needed: every cross-reference of every retained line resolves inside the two delivered spans.  Nothing was omitted from any retained span, so the package carries no omission record.  No retained line contains a citation, so the wrapper carries no bibliography and no bibliography entry.  No retained line contains a figure environment, an image-inclusion command or drawing code, so no image is delivered and the package declares no asset dependency.  Editorial formatting only, in addition: an AMS \texttt{amsart} preamble that restates the article's own declarations and macros used by the retained lines, namely the environments \texttt{lemma} on separate counters, together with the standard packages those lines need.  The retained environments print in this document as Lemma 1 in order of appearance, whereas the article prints the same items with the numbering of its own full document.  The complete native source of the article as distributed in the arXiv e-print for this exact version is bundled unchanged as \texttt{sources/204420/source0.tex}.
\begin{lemma}\label{code}Fix $(i,j)\in R.$  Then
\begin{enumerate}
\item $(i,j)\in T_1(u,v)$ if and only if either $s(i,j)$ begins with a block of the form $+10^n+1$ for some $n\geq 0$ or $s(i,j)$ begins with a block of the form
 $+10^n-1$ for some $n\geq 0$ and $(i+n+1,j+n+1)\in T_2(u,v).$
\item $(i,j)\in T_2(u,v)$ if and only if either $s(i,j)$ begins with a block of the form $-10^n-1$ for some $n\geq 0$ or $s(i,j)$ begins with a block of the form
 $-10^n+1$ for some $n\geq 0$ and $(i+n+1,j+n+1)\in T_1(u,v).$
\end{enumerate}
\end{lemma}

\begin{proof} We prove only (1) as the proof of (2) is symmetric. Let us first assume that $(i,j)\in T_1(u,v).$ Then  $u_j\neq v_i$ and \begin{equation}\label{eq1}u_{j+1}\cdots u_j<_A v_{i+1}\cdots v_i <_D u_{j+1}\cdots u_j.\end{equation}
It follows that $u_j<_A v_i$ and $v_{i+1}\leq _D u_{j+1}.$ This implies that $a_{i,j}=1,$  $d_{i+1,j+1}=0 $ and hence $\delta_{i,j}=+1.$ Thus the  sequence $s(i,j)$ begins in $+1.$ If $s(i,j)$ begins in a block of the form $+10^n+1$ then there is nothing further to prove. So assume that $s(i,j)$ begins in a block of the form $+10^n-1.$ We must show that in this case $(i+n+1,j+n+1)\in T_2(u,v).$ As $s(i,j)$ begins in $+10^n-1$ we have that $\delta_{i,j}=+1,$ $\delta_{i+m,j+m}=0$ for $m=1,\ldots ,n$ and $\delta_{i+n+1,j+n+1}=-1.$  Now
$\delta_{i,j}=+1 \Leftrightarrow a_{i,j}=+1$ and $d_{i+1,j+1}=0$ and hence
\begin{align} \delta_{i,j}=+1 \Leftrightarrow \begin{cases}u_j<_A v_i\\v_{i+1}\leq_D u_{j+1}.\end{cases}\end{align}

Similarly 
\begin{align}\delta_{i+n+1,j+n+1}=-1\Leftrightarrow \begin{cases} v_{i+n+1}\leq_A u_{j+n+1}\\ u_{j+n+2}<_Dv_{i+n+2}.\end{cases}\end{align}
On the other hand, $\delta_{i+m,j+m}=0$ gives rise to two cases depending on whether $a_{i+m,j+m}$ and $d_{i+m,j+m}$ are both equal to $0$ or both equal to $1.$ We find
\begin{align} \delta_{i+m,j+m}=0 \Leftrightarrow \mbox{case (a)}\,\begin{cases}u_{j+m}<_A v_{i+m}\\ u_{j+m+1}<_D v_{i+m+1}\end{cases}\,\,\,\,\mbox{case (b)}\begin{cases}v_{i+m}\leq _A u_{j+m}\\ v_{i+m+1}\leq_D u_{j+m+1}\end{cases}
\end{align}
We begin with $\delta_{i+n+1,j+n+1}=-1.$ Since $u_{j+n+2}<_D v_{i+n+2}$ it follows that $u_{j+n+2}\cdots u_{j+n+1}$ and $v_{i+n+2}\cdots v_{i+n+1}$ begin in different letters and $u_{j+n+2}\cdots u_{j+n+1}<_D v_{i+n+2}\cdots v_{i+n+1}.$ 
To prove that $(i+n+1,j+n+1)\in T_2(u,v)$ it remains to show that $v_{i+n+2}\cdots v_{i+n+1} <_A u_{j+n+2}\cdots u_{j+n+1}.$
If $v_{i+n+1}<_A u_{j+n+1},$ then $v_{i+n+2}\cdots v_{i+n+1}<_A u_{j+n+2}\cdots u_{j+n+1}$ as required. Thus we may assume that
$v_{i+n+1}=u_{j+n+1}.$ If  $v_{i+m}=u_{j+m}$ for each $m=1,\ldots n+1,$  then since $u_{j+n+2}<_Dv_{i+n+2}$ it follows that $(u_{j+1}\cdots u_j)^\omega <_D (v_{i+1}\cdots v_i)^\omega$ and hence $u_{j+1}\cdots u_j <_D v_{i+1}\cdots v_i,$ a contradiction to (\ref{eq1}). 
Thus there exists a largest $1\leq p \leq n$ such that $v_{i+p}\neq  u_{j+p}$ and $v_{i+m}=u_{j+m}$ for each $m=p+1,\ldots ,n+1.$
Then when $\delta_{i+p,j+p}=0,$ we must be in case (b) since in case (a) we have $u_{j+p+1}<_D v_{i+p+1}$ contradicting that $u_{j+p+1}= v_{j+p+1}.$
Thus $ v_{i+p}<_A u_{j+p}.$   
It follows that $(v_{i+n+2}\cdots v_{i+n+1})^{-\omega} <_A (u_{j+n+2}\cdots u_{j+n+1})^{-\omega}$ and hence $v_{i+n+2}\cdots v_{i+n+1}<_A  u_{j+n+2}\cdots u_{j+n+1}$ as required.

We now prove the converse of (1). First assume that $s(i,j)$ begins in $(+1)0^n(+1)$ for some $n\geq 0.$ Then $\delta_{i,j}=1,$ $\delta_{i+m,j+m}=0$ for $m=1,\ldots, n$ and $\delta_{i+n+1,j+n+1}=1.$ As $\delta_{i,j}=1,$  it follows that $u_j<_A v_i$ which implies that $u_{j+1}\cdots u_j$ and $v_{i+1}\cdots v_i$ terminate in distinct letters and  $u_{j+1}\cdots u_j<_A v_{i+1}\cdots v_i.$ To show that $(i,j)\in T_1(u,v)$ it remains to show that 
$v_{i+1}\cdots v_i<_D u_{j+1}\cdots u_j.$   Since $d_{i+1,j+1}=0$ we have that $v_{i+1}\leq _D u_{j+1}.$ If $v_{i+1}<_Du_{j+1}$ then we are done. Otherwise if $v_{i+1}=u_{j+1},$
 then $a_{i+1,j+1}=0$ and hence $d_{i+2,j+2}=0$ (since $\delta_{i+1,j+1}=0).$ So $v_{i+2}\leq_D u_{j+2}.$ If $v_{i+2}<_D u_{j+2}$ then we are done, otherwise if $v_{i+2}=u_{j+2},$ then $a_{i+2,j+2}=0$ and hence $d_{i+3,j+3}=0.$  Continuing in this way, either there exists $m\leq n$ such that $v_{i+k}=u_{j+k}$ for $k=1,\ldots ,m-1$ and $v_{i+m}<_D u_{j+m}$ in which case we have $v[i+1]<_D u[j+1]$ or we have $v_{i+k}=u_{j+k}$ for each $k=1,\ldots ,n.$ In particular, $a_{i+n,j+n}=0$ and hence $d_{i+n+1,j+n+1}=0$ which implies that $v_{i+n+1}\leq_D u_{j+n+1}.$ But since $\delta_{i+n+1,j+n+1}=1,$ it follows that $a_{i+n+1,j+n+1}=1$ which implies in particular that $v_{i+n+1}\neq u_{j+n+1}.$ Hence  $v_{i+n+1}<_D u_{j+n+1}$ as required. 
 
 It remains to show that if $s(i,j)$ begins in $+10^n-1$ and $(i+n+1,j+n+1)\in T_2(u,v),$ then $(i,j)\in T_1(u,v).$ The proof here is essentially symmetric to our earlier proof that if $s(i,j)$ begins in $+10^n-1$  and $(i,j)\in T_1(u,v)$ then $(i+n+1,j+n+1)\in T_2(u,v).$ 
First since $(i+n+1,j+n+1)\in T_2(u,v)$ it follows that $v_{i+n+2}\cdots v_{i+n+1}$ and $u_{j+n+2}\cdots u_{j+n+1}$ begin in different letters and
\begin{equation}\label{con}v_{i+n+2}\cdots v_{i+n+1} <_A  u_{j+n+2}\cdots u_{j+n+1}   <_D v_{i+n+2}\cdots v_{i+n+1}.\end{equation}
This time starting with $\delta_{i,j}=+1,$ we have $u_j<_A v_i$ and $v_{i+1}\leq _D u_{j+1}.$ Thus $u_{j+1}\cdots u_j$ and $v_{i+1}\cdots v_i$ end in different letters and $u_{j+1}\cdots u_j <_A v_{i+1}\cdots v_i.$ To show that $(i,j)\in T_1(u,v)$ it remains to show that $v_{i+1}\cdots v_i <_D  u_{j+1}\cdots u_j.$ If $v_{i+1}<_D u_{j+1}$ then we are done. Thus we can assume hereon that $v_{i+1}=u_{j+1}.$ We consider two cases: Case 1,  $v_{i+m}=u_{j+m}$ for each $m=1,\ldots ,n.$ In this case, when $\delta_{i+m,j+m}=0$ we are in case (b).  In particular $v_{i+n+1}\leq _D u_{j+n+1}.$ We claim $v_{i+n+1} <_D u_{j+n+1}$ and hence that $v_{i+1}\cdots v_i <_D  u_{j+1}\cdots u_j.$ If $v_{i+1}<_D u_{j+1}.$  In fact,  if  $v_{i+n+1}= u_{j+n+1},$ then since $u_j<_A v_i$ we have that $u_{j+n+2}\cdots u_{j+n+1} <_Av_{i+n+2}\cdots v_{i+n+1}$ contradicting (\ref{con}).
Case 2, there exists a least $1<p<n$ such that $v_{i+p}\neq u_{j+p}$ and $v_{i+m}=u_{j+m}$ for $m=1,\ldots ,p-1.$ Thus when $\delta_{i+p-1,j+p-1}=0,$ we must be in case (b) and hence $v_{i+p}<_D u_{j+p}$ and hence $v_{i+1}\cdots v_i<_D u_{j+1}\cdots u_j$ as required. This completes the proof of item (1)  of the lemma. \end{proof}

\end{document}
